%% LyX 1.2 created this file.  For more info, see http://www.lyx.org/.
%% Do not edit unless you really know what you are doing.
\documentclass[oneside,english]{amsart}
\usepackage[T1]{fontenc}
\usepackage[latin1]{inputenc}
\usepackage{graphicx}

\makeatletter

%%%%%%%%%%%%%%%%%%%%%%%%%%%%%% LyX specific LaTeX commands.
\providecommand{\LyX}{L\kern-.1667em\lower.25em\hbox{Y}\kern-.125emX\@}

%%%%%%%%%%%%%%%%%%%%%%%%%%%%%% Textclass specific LaTeX commands.
 \theoremstyle{plain}    
 \newtheorem{thm}{Theorem}[section]
 \numberwithin{equation}{section} %% Comment out for sequentially-numbered
 \numberwithin{figure}{section} %% Comment out for sequentially-numbered
 \theoremstyle{definition}
  \newtheorem{example}[thm]{Example}
 \theoremstyle{definition}
  \newtheorem*{example*}{Example}

%%%%%%%%%%%%%%%%%%%%%%%%%%%%%% User specified LaTeX commands.
\usepackage{epic,eepic,epsfig}
\usepackage{babel}
\makeatother
\begin{document}

\title{Lectures 11 \& 12 --- $M/M/m$ contd, $M/M/m/m$ and $M/G/1$}


\author{Neil Lawrence}

\maketitle
\begin{example*}
Consider the telephone trunking problem. A telephone exchange A is
to server 10,000 subscribers in a nearby exchange as shown in Figure
\ref{fig:exchange}. Suppose that during the busy hour, the subscribers
generate a Poisson traffic with a rate of 10 calls per minute, which
requires trunk lines to exchange B for an average holding time of
3 minutes per call. Determine the number of trunk lines needed for
a grade of service (probability of waiting) of 0.01.

Average call rate is $\lambda =10$ calls/minute. Average holding
time is $\tau =3$ minutes. Offered load $a=\lambda \tau $ is 30
Erlangs.

From the Erlang C curves it looks to be about 44.

%
\begin{figure}
\begin{center}\input{./diagrams/exchange.eepic}\end{center}


\caption{The telephone exchange.\label{fig:exchange}}
\end{figure}

\end{example*}
\begin{example}
A customer service centre has an average holding time for calls of
8 minutes. In the busy hour, calls arrive at the centre at a rate
of one very two minutes. Inter-arrival times and service times are
exponentially distributed, how many personel are required to keep
the probability of waiting at 0.1.

From the plot 8.
\end{example}

\subsection*{Average number in System.}

The average number of calls in the Erlang Delay System

\[
\bar{N}=\frac{\rho }{1-\rho }C(m,a)+m\rho \]
 These terms represent the average queue length and the average number
of busy lines. \[
\bar{N}_{Q}=\frac{\rho }{1-\rho }C(m,a)\]
 So the average waiting time is simply \[
\bar{W}=\frac{\bar{N}_{Q}}{\lambda }=\frac{C(m,a)}{m\mu (1-\rho )}\]


\begin{example*}
So what is the average waiting time for the telephone service centre
example? 
\end{example*}
$m=8$, $\mu =\frac{1}{8}$, $\rho =\frac{\lambda }{m\mu }=\frac{1}{2}$. 

delay is $\frac{0.1}{8\frac{1}{8}(1-\frac{1}{2})}=0.2$ minutes, or
12 seconds. 


\section*{Erlang Loss $M$/$M$/$m$/$m$}

%
\begin{figure}
\begin{center}\includegraphics[  scale=0.7,bb = 0 0 200 100, draft, type=eps]{./diagrams/erlangLoss.eps not found!}\end{center}


\caption{Flow chart of the Erlang Loss System.}
\end{figure}


The Erlang loss system may be analysed as a birth-death process with
birth rate \[
\lambda _{k}=\left\{ \begin{array}{cc}
 \lambda  & k=0,1,\dots ,m-1\\
 0 & k\geq m\end{array}
\right.\]
 and death rate \[
\mu _{k}=\left\{ \begin{array}{cc}
 k\mu  & k=0,1,\dots ,m\\
 0 & k>m\end{array}
\right.\]
 {[}\emph{Note: for the lower line, it doesn't matter what $\mu _{k}$
is.}{]} 

{\tiny }%
\begin{figure}
\begin{center}\input{./diagrams/erlangLoss.eepic}\end{center}


\caption{Markov chain for the Erlang loss system, $M$/$M$/$m$/$m$}
\end{figure}
{\tiny \par}

The offered load, $a=\frac{\lambda }{\mu }$, is again dimensionless
with units of Erlangs. 


\subsection*{Steady State Probabilities}

Using the detailed balance equations we find \begin{eqnarray*}
P(k) & = & \frac{\lambda _{k}}{\mu _{k}}P(k-1)\\
 & = & \frac{\lambda }{k\mu }P(k-1)\\
 & = & \frac{a}{k}P(k-1)
\end{eqnarray*}
 Which means that \begin{eqnarray*}
P(k) & = & \prod _{n=1}^{k}\frac{a}{n}P(0)\\
 & = & \frac{a^{k}}{k!}P(0)
\end{eqnarray*}
 The normalisation condition gives us \[
P(0)\sum _{k=0}^{m}\frac{a^{k}}{k!}=1\]
\[
P(0)=\frac{1}{\sum _{k=0}^{m}\frac{a^{k}}{k!}}\]
 which leads to the truncated Poisson distribution \[
P(k)=\frac{\frac{a^{k}}{k!}}{\sum _{n=0}^{m}\frac{a^{n}}{n!}}\]
 


\subsection*{Performance Measures for Erlang B}

For the blocking probability and the Erlang B curves we are interested
in the probability that when we arrive the queue is full \[
B(m,a)=P(m)=\frac{\frac{a^{m}}{m!}}{\sum _{n=0}^{m}\frac{a^{n}}{n!}}\]
 {\small }%
\begin{figure}
\begin{center}\input{./diagrams/erlangBCurves0to20.eepic}\end{center}


\caption{Erlang B curves.}
\end{figure}
{\small \par}

{\small }%
\begin{figure}
\begin{center}\input{./diagrams/erlangBCurves0to100.eepic}\end{center}


\caption{Erlang B curves.}
\end{figure}
{\small \par}


\subsection*{Carried Load}

The carried load may be found as \[
a^{\prime }=a(1-P(m))=a(1-B(m,a))\]
 The carried load per line is the trunk occupancy, \[
\rho =\frac{a^{\prime }}{m},\]
 or the utilization factor --- it is a measure of efficiency. 

\begin{example*}
Load tripling.

For a blocking probability of 0.01, what is the number of trunks needed
for 6 erlangs of offered load? Answer 12

For a blocking probability of 0.01, what is the number of trunks needed
for 18 erlangs? Answer 30

So load increase of 200 \textbackslash{}\% only lead to a trunk increase
of 150 \textbackslash{}\%. Why? Let's look at the (approximate) utilization
factors. 

\begin{eqnarray*}
\rho _{1} & = & \frac{6}{12}\left(1-0.01\right)=0.495\\
\rho _{2} & = & \frac{18}{30}\left(1-0.01\right)=0.594
\end{eqnarray*}
 The 30 trunk system is thus more efficient.

This explains why it is cheaper to call the U.S.A. than Italy.

$M$/$G$/$1$ Queue
\end{example*}

\subsection*{Pollaczek-Khinchin formula}

If service times are general we cannot use birth death processes to
analyse the system. 

Instead we derive the Pollaczek-Khinchin formula. This formula gives
an expression for the expected waiting time in the queue. \begin{equation}
\bar{W}=\frac{\lambda \left\langle X^{2}\right\rangle }{2(1-\rho )}\end{equation}
 where $\bar{W}$ is expected waiting time, $\left\langle X^{2}\right\rangle =\left\langle X^{2}\right\rangle _{p(X)}$,
$\rho =\lambda \bar{X}$ and $\bar{X}=\left\langle X\right\rangle _{p(X)}$. 


\section*{Proof of Pollaczek-Khinchin Formula}

We will utilise a new concept known as Mean Residual Service Time. 


\subsection*{The Residual Service Time}

Some definitions: 

\begin{tabular}{rcl}
 $W_{i}$&
=&
 Waiting time in queue of $i$th customer.\\
 $R_{i}$&
=&
 `Residual service time' seen by the $i$th customer.\\
 $X_{i}$&
=&
 Service time of the $i$th customer.\\
 $k_{i}$&
=&
 Number of customers in the system when the $i$th customer arrives. \\
\end{tabular}

%
\begin{figure}
\begin{center}\input{diagrams/residualServiceOverview.eepic}\end{center}


\caption{Graphical representation of residual service time.}
\end{figure}


From these definitions we may immediately write down \[
W_{i}=R_{i}+\sum _{j=i-(k_{i}-1)}^{i-1}X_{j}\]
 summing both sides across $i$ we obtain \[
\sum _{i=1}^{I}W_{i}=\sum _{i=1}^{I}R_{i}+\sum _{i=1}^{I}\sum _{j=i-(k_{i}-1)}^{i-1}X_{j}\]
 dividing through by $I$ gives us \[
\sum _{i=1}^{I}\frac{W_{i}}{I}=\sum _{i=1}^{I}\frac{R_{i}}{I}+\sum _{i=1}^{I}\sum _{j=i-(k_{i}-1)}^{i-1}\frac{X_{j}}{I}\]
 in the limit as $I\rightarrow \infty $ this becomes \[
\lim _{I\rightarrow \infty }\sum _{i=1}^{I}\frac{W_{i}}{I}=\lim _{I\rightarrow \infty }\sum _{i=1}^{I}\frac{R_{i}}{I}+\lim _{I\rightarrow \infty }\sum _{i=1}^{I}\frac{(k_{i}-1)X_{i}}{I}\]
 giving \[
\left\langle W\right\rangle =\left\langle R\right\rangle +\left\langle \left(k-1\right)X\right\rangle \]
 Now since the length of the queue as seen by the $i$th arrival is
independent of the $i$th arrival's length this means $\expectation {(k-1)X}=\left\langle k-1\right\rangle \left\langle X\right\rangle $
and since $(k-1)$ is the length of the queue we may write \[
\bar{W}=\bar{R}+\bar{N}_{Q}\bar{X}\]
 Using Little's theorem \[
\bar{N}_{Q}=\lambda \bar{W}\]
 we can obtain \[
\bar{W}=\bar{R}+\lambda \bar{X}\bar{W}\]
 and the fact that $\rho =\lambda \bar{X}$ gives us \[
\bar{W}=\bar{R}+\rho \bar{W}\]
 Leading to \begin{equation}
\bar{W}=\frac{\bar{R}}{1-\rho }\label{eqn:arrivals}\end{equation}
 It remains to calculate $\bar{R}$. 


\subsection*{Residual Service Time from Departures Perspective}

Plot of the residual service time over time. 

%
\begin{figure}

\caption{Graphical computation of mean residual time. \centerline{\input{residualTime.eepic.tex}}}
\end{figure}


The time average of the mean residual service time can be calculated
in terms of the area under the curve, $A(t)$ which in turn is the
area under each of the triangles which may be written \[
A(t)=\frac{1}{2}\sum _{i=1}^{M(t)}X_{i}^{2}\]
 The average service time is the area divided by time, $t$. \begin{eqnarray}
\frac{A(t)}{t} & = & \frac{M(t)}{t}\frac{A(t)}{M(t)}\\
 & = & \frac{1}{2}\frac{M(t)}{t}\frac{\sum _{i=1}^{M(t)}X_{i}^{2}}{M(t)}
\end{eqnarray}
 As $t\rightarrow \infty $ we find $M(t)\rightarrow \infty $ and
\[
\bar{R}=\frac{1}{2}\lambda \left\langle X^{2}\right\rangle .\]
 If we assume that time averages and ensemble averages are equal we
may substitute back into \refeqn{eqn:arrivals} and obtain the Pollaczek-Khinchin
formula \[
\bar{W}=\frac{\lambda \left\langle X^{2}\right\rangle }{2(1-\rho )}\]


This implies that the total waiting time is \[
\bar{T}=\bar{X}+\frac{\lambda \left\langle X^{2}\right\rangle }{2(1-\rho )}\]
 and through Little's formula we can obtain the expected queue length,
$\bar{N}_{Q}$, and the expected number in the system, $\bar{N}$. 


\subsection*{Residual Service Time from Departures Perspective}

Plot of the residual service time over time. 

%
\begin{figure}
\begin{center}\input{./diagrams/residualTime.eepic}\end{center}


\caption{Graphical computation of mean residual time.}
\end{figure}


The time average of the mean residual service time can be calculated
in terms of the area under the curve, $A(t)$ which in turn is the
area under each of the triangles which may be written \[
A(t)=\frac{1}{2}\sum _{i=1}^{M(t)}X_{i}^{2}\]
 The average service time is the area divided by time, $t$. \begin{eqnarray}
\frac{A(t)}{t} & = & \frac{M(t)}{t}\frac{A(t)}{M(t)}\\
 & = & \frac{1}{2}\frac{M(t)}{t}\frac{\sum _{i=1}^{M(t)}X_{i}^{2}}{M(t)}
\end{eqnarray}
 As $t\rightarrow \infty $ we find $M(t)\rightarrow \infty $ and
\[
\bar{R}=\frac{1}{2}\lambda \left\langle X^{2}\right\rangle .\]
 If we assume that time averages and ensemble averages are equal we
may substitute back into eqn (\ref{eqn:arrivals}) and obtain the
Pollaczek-Khinchin formula \[
\bar{W}=\frac{\lambda \left\langle X^{2}\right\rangle }{2(1-\rho )}\]


This implies that the total waiting time is \[
\bar{T}=\bar{X}+\frac{\lambda \left\langle X^{2}\right\rangle }{2(1-\rho )}\]
 and through Little's formula we can obtain the expected queue length,
$\bar{N}_{Q}$, and the expected number in the system, $\bar{N}$. 

\begin{example}
Recall that for an $M$/$M$/$1$system the service times are exponentially
distributed,

\[
P\left(X\right)=\mu \exp \left(-\mu X\right)\]
this means that\[
\bar{X}=\frac{1}{\mu }\]
 then since the variance of the system is $\frac{1}{\mu ^{2}}$\[
\left\langle X\right\rangle ^{2}=\frac{2}{\mu ^{2}}\]
Substituting into the Pollaczek-Khinchin formula we obtain\\
\[
\bar{W}=\frac{\rho }{\mu \left(1-\rho \right)}\]

\end{example}
\begin{example*}
Deterministic service times --- each service takes the same amount
of time\[
P\left(X\right)=\delta \left(X-\frac{1}{\mu }\right)\]
where $\delta \left(\cdot \right)$is the Dirac delta function. For
this system \[
\left\langle X\right\rangle =\frac{1}{\mu }\]
and \[
\left\langle X^{2}\right\rangle =\frac{1}{\mu ^{2}}\]
\emph{i.e.} the variance is zero as we would expect. Substituting
into the Pollaczek-Khinchin formula we obtain\[
\bar{W}=\frac{\rho }{2\mu \left(1-\rho \right)}\]

\end{example*}
In fact we may derive the full range of results for these distributions. 

\begin{tabular}{c|c|c|c}
 Expected value &
 General &
 Exponential &
 Constant\\
\hline
$\bar{W}$&
 $\frac{\lambda }{2(1-\rho )}\expectation {X^{2}}$&
 $\frac{\rho }{\mu (1-\rho )}$&
 $\frac{\rho }{2\mu (1-\rho )}$\\
 $\bar{T}=\bar{W}+\frac{1}{\mu }$&
 $\frac{\lambda }{2(1-\rho )}\expectation {X^{2}}+\frac{1}{\mu }$&
 $\frac{1}{\mu (1-\rho )}$&
 $\frac{2-\rho }{2\mu (1-\rho )}$\\
 $\bar{N_{Q}}=\lambda \bar{W}$&
 $\frac{\lambda ^{2}}{2(1-\rho )}\expectation {X^{2}}$&
$\frac{\rho ^{2}}{(1-\rho )}$&
 $\frac{\rho ^{2}}{2(1-\rho )}$\\
 $\bar{N}=\lambda \bar{W}+\rho =\lambda \bar{T}$&
 $\frac{\lambda ^{2}}{2(1-\rho )}\expectation {X^{2}}+\rho $&
 $\frac{\rho }{1-\rho }$&
 $\frac{\rho (2-\rho )}{2(1-\rho )}$\\
\end{tabular}

\begin{example}
Automatic Repeat Request system --- ARQ System. 

Assume eacdh packet takes 1 millisecond to be transmitted.

The system waits $n-1$ milliseconds for an acknowledgement. If there
is no acknowledgement the system retransmits.

The probability of an error in a packet is taken to be $p$. Errors
are taken to be independent.\\
 Assume that packets arrive as a Poisson process with rate $\lambda $.
\\
 Time interval between transmission and re-transmission is $1+kn$
time units with probability $\left(1-p\right)p^{k}$

The probability of the service time is\[
P\left(X=1+kn\right)=\left(1-p\right)p^{k}\]
 The expected service time is therefore\begin{eqnarray*}
\left\langle X\right\rangle  & = & \sum _{k=0}^{\infty }\left(1+kn\right)\left(1-p\right)p^{k}\\
 & = & \left(1-p\right)\left(\sum _{k=0}^{\infty }p^{k}+n\sum _{k=0}^{\infty }kp^{k}\right)\\
 & = & 1+\frac{np}{1-p}
\end{eqnarray*}
 since $\sum _{k=0}^{\infty }p^{k}=\frac{1}{1-p}$ and $\sum _{k=0}^{\infty }kp^{k}=\frac{p}{\left(1-p\right)^{2}}$.
Thus we can write that\[
\rho =\frac{\lambda \left(1-p\right)}{1-p+np}\]
and thus \[
1-\rho =\frac{1-\lambda -p\left(n-1-\lambda \right)}{1-p+np}\]
The second moment of the service time is\begin{eqnarray*}
\left\langle X^{2}\right\rangle  & = & \sum _{k=0}^{\infty }\left(1+kn\right)^{2}\left(1-p\right)p^{k}\\
 & = & \left(1-p\right)\left(\sum _{k=0}^{\infty }p^{k}+2n\sum _{k=0}^{\infty }kp^{k}+n^{2}\sum _{k=0}^{\infty }k^{2}p^{k}\right)\\
 & = & 1+\frac{2np}{1-p}+\frac{n^{2}\left(p+p^{2}\right)}{\left(1-p\right)^{2}}
\end{eqnarray*}
since \[
\sum _{k=0}^{\infty }k^{2}p^{k}=\frac{p+p^{2}}{\left(1-p\right)^{3}}\]
 Substitiuting into the P-K formula gives a very ugly formula!%
\begin{figure}
\begin{center}\input{./diagrams/ARQ.eepic}\end{center}


\caption{The ARQ System.}
\end{figure}
\end{example}

\end{document}
